\documentclass[11pt]{article}
\usepackage[T1]{fontenc}
\usepackage{lmodern,microtype,amsmath,amssymb,amsthm,mathtools,tikz-cd}
\usepackage[margin=1in]{geometry}
\usepackage{xurl,hyperref}
\hypersetup{colorlinks=true,linkcolor=blue,citecolor=blue,urlcolor=blue,
pdftitle={Syntax and deductively closed theories -- Version 25},pdfauthor={Florian Lengyel}}
\newtheorem{theorem}{Theorem}[section]
\newtheorem{proposition}[theorem]{Proposition}
\newtheorem{lemma}[theorem]{Lemma}
\theoremstyle{definition}\newtheorem{definition}[theorem]{Definition}
\newcommand{\N}{\mathbb N}
\newcommand{\Crec}{\mathcal C_{\mathrm{rec}}}
\newcommand{\Da}{\mathcal D_a}
\newcommand{\Tm}{\mathsf{Tm}}
\newcommand{\Fm}{\mathsf{Fm}}
\newcommand{\FV}{\operatorname{FV}}
\newcommand{\Two}{\mathbf 2}
\newcommand{\Adm}{\mathsf{Adm}}
\newcommand{\im}{\operatorname{im}}
\newcommand{\ran}{\operatorname{ran}}
\newcommand{\Tot}{\operatorname{Tot}}
\newcommand{\Split}{\operatorname{Split}}
\newcommand{\Par}{\operatorname{Par}}
\newcommand{\Mono}{\operatorname{Mono}}
\newcommand{\List}{\mathsf L}
\newcommand{\nil}{\mathsf{nil}}
\newcommand{\cons}{\mathsf{cons}}
\newcommand{\mem}{\mathsf{mem}}
\newcommand{\del}{\mathsf{del}}
\newcommand{\ok}{\mathsf{ok}}
\newcommand{\fail}{\mathsf{fail}}
\begin{document}
\appendix
\section{Syntax and deductively closed theories}
\label{app:syntax-theories}
\noindent\textit{Version 25, 1 November 2026.}\par\medskip

We define first-order syntax by categorical structure and construct
the theory generated by an axiom subobject. The construction uses
parameterized initial algebras for syntax and lists, an initial algebra
of derivations, and regular images.

\subsection{Categorical hypotheses and notation}
\label{app:categorical-hypotheses}

Let $\mathcal D$ be a regular category with finite coproducts.
Thus $\mathcal D$ has finite limits, every morphism factors as a
regular epimorphism followed by a monomorphism, and regular
epimorphisms are stable under pullback. A regular epimorphism is a
coequalizer; the monic part of the factorization is the image.
For the construction of first-order syntax, assume also that finite
coproducts are disjoint and stable under pullback. In particular,
a morphism into $X+Y$ decomposes its source as the coproduct of its
two inverse images. These additional coproduct hypotheses are not
needed for Theorem~\ref{thm:least-theorem-subobject}.

Composition is $gf=g\circ f$. Finite products and coproducts are
chosen and parenthesized to the left. Rebracketing uses the canonical
associativity isomorphisms determined by their projections or injections;
terminal factors are removed by the canonical product unit isomorphisms.
For $X\in\mathcal D$, write $X^0=1$, $X^1=X$, and
$X^{n+1}=X^n\times X$ for $n\geq1$.
For $f:X\to Y$, the morphism $f^n:X^n\to Y^n$ is the corresponding
product morphism, with $f^0=1_1$ and $f^1=f$.
Thus powers of objects denote finite products.

Put $\Two=1+1$, with injections
$\mathbf{f},\mathbf{t}:1\to\Two$. Pulling this coproduct back along
$\pi_1:\Two\times\Two\to\Two$ decomposes $\Two\times\Two$ into two copies of
$\Two$. Decomposing each copy as $1+1$ gives the isomorphism
\[
 \gamma=[\langle\mathbf{f},\mathbf{f}\rangle,
          \langle\mathbf{f},\mathbf{t}\rangle,
          \langle\mathbf{t},\mathbf{f}\rangle,
          \langle\mathbf{t},\mathbf{t}\rangle]:
          1+1+1+1\xrightarrow{\cong}\Two\times\Two.
\]
Define the Boolean morphisms by
\[
 \neg_2=[\mathbf{t},\mathbf{f}]:\Two\to\Two,\qquad
 \wedge_2=[\mathbf{f},\mathbf{f},\mathbf{f},\mathbf{t}]\gamma^{-1},
 \qquad
 \vee_2=[\mathbf{f},\mathbf{t},\mathbf{t},\mathbf{t}]\gamma^{-1},
\]
where $\wedge_2,\vee_2:\Two\times\Two\to\Two$.
For $\mathbf{a},\mathbf{b}$ each equal to one of the injections
$\mathbf{f},\mathbf{t}$, let $j_{\mathbf{a}\mathbf{b}}:1\to1+1+1+1$
denote the corresponding summand injection. The definition of $\gamma$ gives
$\gamma j_{\mathbf{a}\mathbf{b}}=\langle\mathbf{a},\mathbf{b}\rangle$.
Consequently,
\[
 \wedge_2\langle\mathbf{a},\mathbf{b}\rangle
 =[\mathbf{f},\mathbf{f},\mathbf{f},\mathbf{t}]
    j_{\mathbf{a}\mathbf{b}},\qquad
 \vee_2\langle\mathbf{a},\mathbf{b}\rangle
 =[\mathbf{f},\mathbf{t},\mathbf{t},\mathbf{t}]
    j_{\mathbf{a}\mathbf{b}}.
\]
The defining equations of a copairing therefore give
\[
\begin{array}{cc|cc}
 \mathbf{a}&\mathbf{b}&
 \wedge_2\langle\mathbf{a},\mathbf{b}\rangle&
 \vee_2\langle\mathbf{a},\mathbf{b}\rangle\\\hline
 \mathbf{f}&\mathbf{f}&\mathbf{f}&\mathbf{f}\\
 \mathbf{f}&\mathbf{t}&\mathbf{f}&\mathbf{t}\\
 \mathbf{t}&\mathbf{f}&\mathbf{f}&\mathbf{t}\\
 \mathbf{t}&\mathbf{t}&\mathbf{t}&\mathbf{t}.
\end{array}
\]
Also $\neg_2\mathbf{f}=\mathbf{t}$ and
$\neg_2\mathbf{t}=\mathbf{f}$.

These morphisms make $\Two$ a Boolean algebra object in $\mathcal D$.
Indeed, distributivity expresses each finite power $\Two^n$ as a
coproduct of $2^n$ copies of $1$, indexed by the ordered choices of
$\mathbf{f}$ or $\mathbf{t}$ in each coordinate. The resulting
injections are jointly epimorphic. The equations above verify the
commutativity, idempotence, unit, absorption, and complement identities
after precomposition with the injections for their respective sources.
For associativity of $\wedge_2$, both composites on $\Two^3$ are
$\mathbf{t}$ on the summand labelled $(\mathbf{t},\mathbf{t},\mathbf{t})$
and $\mathbf{f}$ on the other seven summands. For associativity of
$\vee_2$, both composites are $\mathbf{f}$ on the summand labelled
$(\mathbf{f},\mathbf{f},\mathbf{f})$ and $\mathbf{t}$ on the others.
For distributivity of $\wedge_2$ over $\vee_2$, both composites are
$\mathbf{t}$ on the summands labelled
$(\mathbf{t},\mathbf{f},\mathbf{t})$,
$(\mathbf{t},\mathbf{t},\mathbf{f})$, and
$(\mathbf{t},\mathbf{t},\mathbf{t})$; both are $\mathbf{f}$ on the others.
For distributivity of $\vee_2$ over $\wedge_2$, both composites are
$\mathbf{f}$ on the summands labelled
$(\mathbf{f},\mathbf{f},\mathbf{f})$,
$(\mathbf{f},\mathbf{f},\mathbf{t})$, and
$(\mathbf{f},\mathbf{t},\mathbf{f})$; both are $\mathbf{t}$ on the others.
Joint epimorphy then gives these identities as equalities of morphisms
in $\mathcal D$.

Let $p:X\to\Two$. Write $[p=\mathbf{t}]\hookrightarrow X$ for its
pullback along $\mathbf{t}$; write $[p=\mathbf{f}]\hookrightarrow X$ for its
pullback along $\mathbf{f}$.
Denote these inclusions by $i_{\mathbf{t}}$ and $i_{\mathbf{f}}$.
The coproduct hypotheses give the isomorphism
\[
 \delta_p=[i_{\mathbf{t}},i_{\mathbf{f}}]:
 [p=\mathbf{t}]+[p=\mathbf{f}]\xrightarrow{\cong}X.
\]
For branches $f,g:X\to Y$, define the conditional by
\[
 \mathsf{if}(p;f,g)=
 [f i_{\mathbf{t}},g i_{\mathbf{f}}]\delta_p^{-1}:X\to Y.
\]
It is the unique morphism $h:X\to Y$ satisfying
$h i_{\mathbf{t}}=f i_{\mathbf{t}}$ and
$h i_{\mathbf{f}}=g i_{\mathbf{f}}$.

We use variables to abbreviate projections from a displayed product.
An expression built from them denotes a composite of morphisms on
that product. An equation between such expressions is an equality
of morphisms with that common source. A morphism $c:1\to Y$ used
on a source $X$ denotes $c!_X:X\to Y$.

Fix an object $V\in\operatorname{Ob}(\mathcal D)$ and a morphism
$\epsilon:V\times V\to\Two$ in $\mathcal D$ such that the square
\[
\begin{tikzcd}
 V\arrow[r,"{!_V}"]\arrow[d,hook,"\Delta_V"'] &
 1\arrow[d,hook,"\mathbf{t}"]\\
 V\times V\arrow[r,"\epsilon"'] & \Two
\end{tikzcd}
\]
is a pullback. Thus $V$ is equipped with decidable equality:
its diagonal is a complemented subobject of $V\times V$, with
complement given by the pullback of $\mathbf{f}:1\to\Two$
along $\epsilon$.
For $0\leq k\leq d$, fix objects
$\Sigma_k,\Pi_k\in\operatorname{Ob}(\mathcal D)$ to index the
$k$-ary function and relation constructors, respectively.
Here $d$ is a fixed nonnegative integer. Equality is a logical symbol.

For an endofunctor on the $r$-fold product category $\mathcal D^r$
defined using finite products, coproducts, and specified constant objects, its version on
$(\mathcal D/P)^r$ is obtained
by replacing each constant $X$ by $P\times X\to P$ and using
products and coproducts in that slice. An initial algebra is
\emph{stable under parameters} if, for every $P\in\mathcal D$,
its pullback along $P\to1$ is initial for this slice endofunctor.
For list objects, this is the parameterized recursion property studied by
Cockett~\cite{CockettLists}; its formulation by pullback into slices is
recalled by Maietti~\cite[Definition 2.4 and the following paragraph]{MaiettiLists}.

\subsection{Terms, formulas, and free variables}
\label{app:first-order-syntax}

Initial many-sorted algebras are used to specify abstract syntax by
Goguen, Thatcher, Wagner, and Wright~\cite[Section 2]{GTWW}.
Goguen and Burstall~\cite[Section 4]{GoguenBurstall} apply this method
to the term and sentence sorts of logical languages.
Here the signature objects, syntax objects, and constructor morphisms
lie in $\mathcal D$.

Initial-algebra accounts of binding and substitution are developed by
Fiore, Plotkin, and Turi~\cite[Sections 2--4]{FPT} and by
Gabbay and Pitts~\cite[Section 6]{GabbayPitts}.
Katsumata~\cite[Section 2, Example 1(3)]{Katsumata} gives a first-order
instance with a quantifier. These treatments identify expressions up to
renaming of bound variables. The syntax defined below retains variable
names, and formula substitution is defined on an admissibility subobject.

Define $K:\mathcal D^2\to\mathcal D^2$ by
$K(X,Y)=\bigl(K_1(X,Y),K_2(X,Y)\bigr)$, where
\begin{align}
 K_1(X,Y)&=V+\coprod_{k=0}^{d}\Sigma_k\times X^k,
                                                        \label{eq:term-functor}\\
 K_2(X,Y)&=X^2+\coprod_{k=0}^{d}\Pi_k\times X^k
                  +Y+Y^2+V\times Y.                    \label{eq:formula-functor}
\end{align}
On a morphism $(f,g):(X,Y)\to(X',Y')$, use
$1_V$, $1_{\Sigma_k}\times f^k$, $f^2$,
$1_{\Pi_k}\times f^k$, $g$, $g^2$, and $1_V\times g$ on
the respective summands. Product and coproduct laws give the
functor laws.

\begin{definition}\label{def:categorical-syntax}
A raw named-variable syntax object is an initial $K$-algebra
$\iota:K(\Tm,\Fm)\to(\Tm,\Fm)$ stable under parameters.
We also require a list object $W=\List V$, given by an initial
algebra $[\nil,\cons]:1+V\times W\to W$ stable under parameters.
\end{definition}

The syntax constructors are the following morphisms in $\mathcal D$:
\begin{gather*}
 \eta:V\to\Tm,\qquad f_k:\Sigma_k\times\Tm^k\to\Tm,\\
 {\mathsf{eq}}:\Tm^2\to\Fm,\qquad
 r_k:\Pi_k\times\Tm^k\to\Fm,\\
 \neg:\Fm\to\Fm,\qquad
 {\Rightarrow}:\Fm^2\to\Fm,\qquad
 {\forall}:V\times\Fm\to\Fm.
\end{gather*}
The morphism $\eta$ is the variable constructor.
The quantifier constructor $\forall$ is indexed
by $V$; its binding behaviour is specified by the free-variable and
substitution morphisms constructed below.
For every $K$-algebra $\alpha:K(X,Y)\to(X,Y)$ there is a unique
$(u,v):(\Tm,\Fm)\to(X,Y)$ making the square
\[
\begin{tikzcd}[column sep=large]
 K(\Tm,\Fm)\arrow[r,"\iota"]\arrow[d,"{K(u,v)}"'] &
 (\Tm,\Fm)\arrow[d,"{(u,v)}"]\\
 K(X,Y)\arrow[r,"\alpha"'] & (X,Y)
\end{tikzcd}
\]
commute. The same property holds in every slice with parameters.
Induction means applying this property to a monic subalgebra:
the unique morphism into that subalgebra is an inverse to its
inclusion. Thus parameterized induction is also available.

List initiality with parameters gives concatenation
$+\!+:W^2\to W$, membership $\mem:V\times W\to\Two$, and
deletion $\del:V\times W\to W$ by
\begin{align*}
 \nil+\!+w&=w,&
 \cons(y,w)+\!+w'&=\cons(y,w+\!+w'),\\
 \mem(x,\nil)&=\mathbf{f},&
 \mem(x,\cons(y,w))&=\epsilon(x,y)\vee_2\mem(x,w),\\
 \del(x,\nil)&=\nil,&
 \del(x,\cons(y,w))&=
 \mathsf{if}\bigl(\epsilon(x,y);\del(x,w),
                         \cons(y,\del(x,w))\bigr).
\end{align*}
In the nil equations the unused variables specify the source of the
constant morphism. These are the algebra-morphism equations of the
corresponding folds. Write $[x]=\cons(x,\nil)$.
Parameterized list induction gives the identities
\begin{align*}
 \mem(x,w+\!+w')&=\mem(x,w)\vee_2\mem(x,w'),\\
 \mem(x,\del(y,w))&=\neg_2\epsilon(x,y)\wedge_2\mem(x,w).
\end{align*}

Give $(W,W)$ the $K$-algebra defined by singleton lists at variables,
ordered concatenation at function and relation symbols and equality,
the identity at negation, concatenation at implication, and deletion
at quantification. Its unique algebra morphism defines
$\FV_{\Tm}:\Tm\to W$ and $\FV_{\Fm}:\Fm\to W$. Explicitly,
\begin{align*}
 \FV_{\Tm}(\eta x)&=[x],&
 \FV_{\Tm}(f_k(s,\vec t))&=\FV_{\Tm}(t_1)+\!+\cdots+\!+\FV_{\Tm}(t_k),\\
 \FV_{\Fm}(\mathsf{eq}(u,t))&=\FV_{\Tm}(u)+\!+\FV_{\Tm}(t),&
 \FV_{\Fm}(r_k(s,\vec t))&=\FV_{\Tm}(t_1)+\!+\cdots+\!+\FV_{\Tm}(t_k),\\
 \FV_{\Fm}(\neg\varphi)&=\FV_{\Fm}(\varphi),&
 \FV_{\Fm}(\varphi\Rightarrow\psi)&=\FV_{\Fm}(\varphi)+\!+\FV_{\Fm}(\psi),\\
 \FV_{\Fm}(\forall x\,\varphi)&=\del(x,\FV_{\Fm}(\varphi)).&&
\end{align*}
An empty concatenation is $\nil$. These are lists of free-variable
occurrences; repetitions are retained.

\begin{definition}\label{def:sentence-object}
The sentence object $s:S\hookrightarrow\Fm$ is the pullback
\[
\begin{tikzcd}
 S\arrow[r,hook,"s"]\arrow[d] & \Fm\arrow[d,"\FV_{\Fm}"]\\
 1\arrow[r,"\nil"'] & W.
\end{tikzcd}
\]
\end{definition}
Negation restricts to $\nu:S\to S$ by the displayed free-variable
equation. Existential quantification is the derived morphism
$\exists=\neg\circ\forall\circ(1_V\times\neg):V\times\Fm\to\Fm$.

\subsection{Admissible substitution}
\label{app:substitution}

We construct substitution of one term for one variable, as required
by the logical schemes below. Put $E=V\times\Tm$, with projections
$x:E\to V$ and $t:E\to\Tm$. For an object $Y$, write
$O(Y)=1+Y$, with injections $\fail_Y:1\to O(Y)$ and
$\ok_Y:Y\to O(Y)$. A morphism $g:Y\to Z$ induces
$O(g)=1_1+g$. For a binary morphism $h:Y^2\to Y$, let
$\widehat h:O(Y)^2\to O(Y)$ return $\ok_Yh$ on the
$Y\times Y$ summand and $\fail_Y$ on the other three summands.
Also let $\widehat\forall:V\times O(\Fm)\to O(\Fm)$ have
restrictions $\fail_{\Fm}$ and $\ok_{\Fm}\forall$.
Distributivity of products over coproducts defines these morphisms.

In $(\mathcal D/E)^2$, use the carriers
\[
 X_T=E\times(\Tm\times\Tm)\longrightarrow E,
 \qquad
 X_F=E\times(\Fm\times O(\Fm))\longrightarrow E.
\]
The first component of each pair records the original expression.
Retaining this component obtains structural recursion from initial-algebra
iteration; compare Gabbay and Pitts~\cite[proof of Corollary 6.7]{GabbayPitts}.
Define a $K_E$-algebra as follows, retaining the same parameter
$(x,t)$ in every operation. At a variable $v$, the pair is
\[
 (\eta v,\mathsf{if}(\epsilon(v,x);t,\eta v)).
\]
At a function symbol, apply its constructor separately to the
original term components and the substituted term components.
At equality and relation symbols, reconstruct the original formula
and apply $\ok_{\Fm}$ to the formula formed from the substituted
term components. At negation and implication, reconstruct the
original formula and apply $O(\neg)$ or
$\widehat{\Rightarrow}$ to the result components.

For quantification, the operation has source
$E\times V\times\Fm\times O(\Fm)$.
Denote its projections by $(x,t,y,\varphi,z)$, with $(x,t)$ the
$E$-coordinate. Define
\[
 C(y,\varphi,x,t)=
 \neg_2\mem(x,\FV_{\Fm}(\varphi))
 \vee_2\neg_2\mem(y,\FV_{\Tm}(t)).
\]
The original component is $\forall y\,\varphi$; the result component is
\begin{equation}\label{eq:quantifier-fold}
 \mathsf{if}\left(\epsilon(y,x);
   \ok_{\Fm}(\forall y\,\varphi),
   \mathsf{if}\bigl(C(y,\varphi,x,t);
          \widehat\forall(y,z),\fail_{\Fm}\bigr)\right).
\end{equation}
In particular, the branch $y=x$ uses the original body and discards
its substitution result.

\begin{proposition}\label{prop:categorical-substitution}
Parameterized initiality determines morphisms
\[
 a_T:\Tm\times E\longrightarrow\Tm,
 \qquad b_F:\Fm\times E\longrightarrow O(\Fm).
\]
Admissibility and formula substitution are defined by the pullback
\[
\begin{tikzcd}
 \Adm\arrow[r,"a_F"]\arrow[d,hook,"j"'] &
 \Fm\arrow[d,hook,"\ok_{\Fm}"]\\
 \Fm\times E\arrow[r,"b_F"'] & O(\Fm).
\end{tikzcd}
\]
The monomorphism $j$ is complemented. These morphisms are uniquely
determined by the preceding algebra and its initial morphism.
\end{proposition}
\begin{proof}
The unique $K_E$-algebra morphism into $(X_T,X_F)$ records an original
and a result component. Composing this algebra morphism with the
projections to the original components gives an endomorphism of the
initial $K_E$-algebra, hence the identity. The result components,
after interchanging the parameter
and syntax factors, are $a_T$ and $b_F$. The pullback defines
$j$ and $a_F$; the coproduct hypotheses make $j$ complemented.
\end{proof}

Put $B=V\times\Fm\times E$ and define
\[
 q=\langle\forall\pi_{V\times\Fm},\pi_E\rangle:B\to\Fm\times E,
 \qquad
 p=\langle\pi_{\Fm},\pi_E\rangle:B\to\Fm\times E.
\]
Let $D_=\hookrightarrow B$ and $D_{\ne}\hookrightarrow B$ be the
two fibers of $\epsilon(y,x)$, and let $D_C=[C=\mathbf{t}]$.
Equation~\eqref{eq:quantifier-fold} gives
\begin{equation}\label{eq:admissibility-subobjects}
 q^*\Adm=D_=\;\vee\;
           (D_{\ne}\wedge D_C\wedge p^*\Adm).
\end{equation}
Here meet is pullback intersection and the displayed disjoint join
is the coproduct inclusion into $B$. On $D_=$ the substitution
result is $\forall y\,\varphi$; on the other summand it is
$\forall(y,a_F\widetilde p)$, where
$\widetilde p$ is the induced morphism to $\Adm$.
Thus quantification binds its named variable and prevents capture
of free occurrences of the replacement term.

We write $\varphi[x:=t]$ only on $\Adm$, and $u[x:=t]$ for
$a_T(u,x,t)$. Define
\[
 e_T=\langle\pi_{\Tm},\pi_V,\eta\pi_V\rangle:
          \Tm\times V\to\Tm\times E,
 \qquad
 e_F=\langle\pi_{\Fm},\pi_V,\eta\pi_V\rangle:
          \Fm\times V\to\Fm\times E.
\]
Identity substitution satisfies
\[
 a_Te_T=\pi_{\Tm},\qquad b_Fe_F=\ok_{\Fm}\pi_{\Fm}.
\]
The second equality also states admissibility. For substitution at
an absent variable, use the subobjects
\begin{align*}
 N_T&=[\mem(x,\FV_{\Tm}(u))=\mathbf{f}]
                            \hookrightarrow\Tm\times E,\\
 N_F&=[\mem(x,\FV_{\Fm}(\varphi))=\mathbf{f}]
                            \hookrightarrow\Fm\times E.
\end{align*}
On $N_T$, $a_T$ is the term projection. On $N_F$, $b_F$ is
$\ok_{\Fm}$ composed with the formula projection, so substitution
is admissible and fixes the formula.

These term and formula statements follow simultaneously by
parameterized induction. Atomic steps use the term equations;
connective steps use the lifted operations. At a quantifier with
equal variable names, \eqref{eq:quantifier-fold} returns the original
formula. For identity substitution with distinct names,
$\FV_{\Tm}(\eta x)=[x]$ makes the capture test true.
For an absent variable with distinct names, the deletion identity
transfers the absence condition to the body. In both cases the
induction hypothesis gives the required successful second branch.

\subsection{Inference rules and deductively closed subobjects}
\label{app:closed-subobjects}

Aczel~\cite[Propositions 1.1.4 and 1.4.5]{Aczel} describes finitary
inductive closure by the conclusions of finite proofs. We express the
rule presentation in a slice category.
Let $\mathcal D$ be regular with finite coproducts, and fix an object
$F$. For an integer $N\geq0$, let
\[
 p_n:R_n\longrightarrow F^n,\qquad
 q_n:R_n\longrightarrow F\quad(0\leq n\leq N)
\]
specify rule instances, their ordered premises, and their conclusions.
The nullary rule instances specify axioms.
For an object $q:Q\to F$ of the slice category $\mathcal D/F$,
form the pullback
\[
\begin{tikzcd}
 B_n(Q,q)\arrow[r,"{b_{n,q}}"]\arrow[d,"{a_{n,q}}"'] &
 Q^n\arrow[d,"q^n"]\\
 R_n\arrow[r,"p_n"'] & F^n.
\end{tikzcd}
\]
Put
\[
 H(Q,q)=\left(\coprod_{n=0}^N B_n(Q,q),\kappa_q\right),\qquad
 \kappa_q=[q_na_{n,q}]_{n=0}^N.
\]
For a slice morphism $h:(Q,q)\to(Q',q')$, define $B_n(h)$ by
\[
 a_{n,q'}B_n(h)=a_{n,q},\qquad
 b_{n,q'}B_n(h)=h^n b_{n,q}.
\]
Pullback uniqueness proves the identity and composition laws, so
$H(h)=\coprod_nB_n(h)$ defines an endofunctor of $\mathcal D/F$.

\begin{definition}\label{def:deductively-closed}
A subobject $m:T\hookrightarrow F$ is closed under these rules if
there is a morphism $\alpha:\coprod_nB_n(T,m)\to T$ with
$m\alpha=\kappa_m$. Equivalently, $(T,m)$ admits an $H$-algebra
structure in $\mathcal D/F$. When $F$ and the rule morphisms specify
a logical calculus, such a subobject is a deductively closed theory
of formulas for that calculus.
\end{definition}

The algebra structure is unique because $m$ is monic. Its nullary
part includes all axioms, and its other parts express closure under
the inference rules. Inclusion of closed subobjects is automatically
an algebra morphism, by composing with the target monomorphism.

The following proof transfers the initial derivation algebra to its
regular image. This is an instance of lifting a factorization system to
algebras for an endofunctor preserving its left class; see, for example,
Wi\ss{}mann~\cite[Lemma 3.6]{Wissmann}.

\begin{theorem}\label{thm:least-theorem-subobject}
With the rule data above, suppose
$\delta:H(P,c)\to(P,c)$ is an initial algebra in $\mathcal D/F$.
The image of its conclusion morphism $c:P\to F$ is the least
subobject of $F$ closed under the rules.
\end{theorem}
\begin{proof}
Factor $c=me$ with $e:P\twoheadrightarrow T$ regular epic and
$m:T\hookrightarrow F$ monic. The endofunctor $H$ preserves
regular epimorphisms: finite products preserve them in a regular
category, each $B_n(e)$ is a pullback of $e^n$, and finite coproducts
preserve coequalizers. Thus $H(e)$ is regular epic.

Writing $|H(P,c)|$ for the underlying object of the slice object,
the equations $c\delta=\kappa_c$ and
$\kappa_mH(e)=\kappa_c$ give the commuting square
\[
\begin{tikzcd}[column sep=large]
 \lvert H(P,c)\rvert\arrow[r,"e\delta"]\arrow[d,two heads,"H(e)"'] &
 T\arrow[d,hook,"m"]\\
 \lvert H(T,m)\rvert\arrow[r,"\kappa_m"']\arrow[ur,dashed,"\alpha"] & F.
\end{tikzcd}
\]
Since $m$ is monic, $e\delta$ coequalizes the kernel pair of $H(e)$.
The regular epimorphism $H(e)$ is its kernel-pair coequalizer, so
there is a unique diagonal $\alpha$. In particular, $m\alpha=\kappa_m$,
so the image is closed. If $u:U\hookrightarrow F$ is another
closed subobject, initiality gives an algebra morphism
$v:(P,c)\to(U,u)$, hence $uv=c$. The image factorization gives
a unique $w:T\to U$ with $uw=m$. This is the required inclusion.
\end{proof}

\subsection{First-order theories}
\label{app:first-order-derivations}

Continue with the syntax structure of
Definition~\ref{def:categorical-syntax} and its constructed substitution.
Put $F=\Fm$. The propositional axiom schemes are those of
H\'ajek and Pudl\'ak~\cite[Section 0.10]{HP}:
\begin{align}
 &\varphi\Rightarrow(\psi\Rightarrow\varphi),
                                                        \label{eq:hilbert-one}\\
 &(\varphi\Rightarrow(\psi\Rightarrow\chi))
   \Rightarrow((\varphi\Rightarrow\psi)
   \Rightarrow(\varphi\Rightarrow\chi)),
                                                        \label{eq:hilbert-two}\\
 &(\neg\psi\Rightarrow\neg\varphi)
   \Rightarrow(\varphi\Rightarrow\psi).
                                                        \label{eq:hilbert-three}
\end{align}
The named-variable quantifier axioms follow
Hirst and Hirst~\cite[Sections 2.11--2.12]{HH}; equality is expressed
by reflexivity and substitution of equal terms in formulas
\cite[Section 3.1]{HH}:
\begin{align}
 &(\forall x\,\varphi)\Rightarrow\varphi[x:=t],
                                                        \label{eq:hilbert-instantiation}\\
 &\forall x(\varphi\Rightarrow\psi)
       \Rightarrow(\varphi\Rightarrow\forall x\,\psi),
                                                        \label{eq:hilbert-distribution}\\
 &t=t,                                                   \label{eq:hilbert-reflexivity}\\
 &(u=t)\Rightarrow
       (\varphi[x:=u]\Rightarrow\varphi[x:=t]).           \label{eq:hilbert-equality}
\end{align}
Each scheme denotes a morphism from its parameter object to $F$.
For the first three these objects are $F^2,F^3,F^2$.
For \eqref{eq:hilbert-instantiation}, use the pullback of
$j:\Adm\hookrightarrow F\times V\times\Tm$ along the
permutation $V\times F\times\Tm\to F\times V\times\Tm$.
For \eqref{eq:hilbert-distribution}, use
$[\mem(x,\FV_{\Fm}(\varphi))=\mathbf{f}]
\hookrightarrow V\times F^2$.
For \eqref{eq:hilbert-reflexivity}, use $\Tm$.
For \eqref{eq:hilbert-equality}, intersect the pullbacks of $j$
along
\[
 (x,\varphi,u,t)\longmapsto(\varphi,x,u),\qquad
 (x,\varphi,u,t)\longmapsto(\varphi,x,t)
\]
from $V\times F\times\Tm^2$ to $F\times E$.
Let $L$ be the coproduct of these seven parameter objects and
$\lambda:L\to F$ their copairing.

Fix an axiom subobject $i_A:A\hookrightarrow S$. Define rule objects
and their premise and conclusion morphisms by
\begin{gather*}
 R_0=L+A,\quad p_0:R_0\to1,\quad q_0=[\lambda,si_A],\\
 R_1=V\times F,\quad p_1=\pi_F:R_1\to F,
             \quad q_1=\forall:R_1\to F,\\
 R_2=F^2,\quad p_2=\langle\pi_1,\Rightarrow\rangle:R_2\to F^2,
             \quad q_2=\pi_2:R_2\to F.
\end{gather*}
These are axiom introduction, generalization, and modus ponens.
All nonlogical axioms are sentences. Let $H_A:\mathcal D/F\to
\mathcal D/F$ be the resulting inference endofunctor.

\begin{theorem}\label{thm:first-order-theory}
Suppose $H_A$ has an initial algebra with conclusion morphism
$c_A:P_A\to F$. Its image $T_A\hookrightarrow F$ is the least
subobject containing the logical and nonlogical axioms and closed
under generalization and modus ponens. The image of the pullback
\[
\begin{tikzcd}
 P_{A,S}\arrow[r]\arrow[d,"c_{A,S}"'] & P_A\arrow[d,"c_A"]\\
 S\arrow[r,hook,"s"'] & F
\end{tikzcd}
\]
is $T_{A,S}=S\times_F T_A\hookrightarrow S$.
\end{theorem}
\begin{proof}
Theorem~\ref{thm:least-theorem-subobject} gives the first assertion.
Regular images are stable under pullback, giving the second.
\end{proof}

Suppose these initial derivation algebras exist for every axiom
subobject of $S$. A \emph{sentence theory} is a subobject
$u:U\hookrightarrow S$ through which $c_{U,S}$ factors.
Then $T_{A,S}$ is the least sentence theory containing $A$.
Indeed, axiom introduction gives $A\subseteq T_{A,S}$; since $T_A$
is closed and contains $T_{A,S}$, leastness gives
$T_{T_{A,S}}\subseteq T_A$. Pullback proves closure of $T_{A,S}$.
If $U$ is a sentence theory containing $A$, leastness gives
$T_A\subseteq T_U$, hence
$T_{A,S}\subseteq T_{U,S}\subseteq U$.
It also follows that $A\subseteq A'$ implies
$T_{A,S}\subseteq T_{A',S}$.

For a restriction category $\mathcal R$, its restriction splitting
$\Split(\mathcal R)$ has objects $(X,e)$, where $e=\bar e$, and
morphisms $f:(X,e)\rightharpoonup(Y,d)$ satisfying $f=dfe$.
Its identity on $(X,e)$ is $e$; such a morphism is total when
$\bar f=e$. Write $\Tot(\Split(\mathcal R))$ for the subcategory
of total morphisms~\cite[Section 2.3.3]{CockettLack}.

In a restriction category $\mathcal S$, restriction idempotents
$a,b:Y\rightharpoonup Y$ are ordered by $a\leq b$ when $ba=a$.
A restriction idempotent $a$ receives $f:X\rightharpoonup Y$ if $af=f$.
The \emph{range} $\ran(f)$ is the least restriction idempotent receiving
$f$. We use ranges satisfying the agreement law
\[
 uf=vf\quad\Longrightarrow\quad u\ran(f)=v\ran(f)
 \qquad(u,v:Y\rightharpoonup Z).
\]
A total morphism $m:I\to Y$ is a \emph{restriction monomorphism}
if it has a partial inverse $r:Y\rightharpoonup I$ with
$rm=1_I$ and $mr=\bar r$.

\begin{lemma}[Regular images and ranges]\label{lem:range-regular-image}
Let $\mathcal S$ be a restriction category with ranges in which every
restriction idempotent splits. Suppose $\mathcal D=\Tot(\mathcal S)$
is regular, and let
\[
 f:X\longrightarrow Y,\qquad
 f=me,\qquad e:X\twoheadrightarrow I,\qquad m:I\hookrightarrow Y
\]
be a regular-image factorization in $\mathcal D$.
If $m$ has a partial inverse $r:Y\rightharpoonup I$ in $\mathcal S$,
then $\ran(f)=mr=\bar r$.
\end{lemma}
\begin{proof}
The restriction idempotent $mr$ receives $f$, since $mrf=mrme=f$.
Let $a:Y\rightharpoonup Y$ be any restriction idempotent receiving $f$.
Choose a splitting $a=js$, where $j:J\to Y$ is total,
$s:Y\rightharpoonup J$, and $sj=1_J$.
Then $f=j(sf)$, and $sf:X\to J$ is total, since
$\overline{sf}=\overline{jsf}=\bar f=1_X$.
The morphism $j$ is monic in $\mathcal D$, so the universal property
of the regular image gives a unique $u:I\to J$ with $ju=m$.
Consequently,
\[
 am=jsju=ju=m,\qquad a(mr)=mr.
\]
Thus $mr\leq a$, proving that $mr$ is the least restriction
idempotent receiving $f$.
\end{proof}

The relation between split ranges and factorizations through restriction
monomorphisms is developed in Cockett, Guo, and
Hofstra~\cite[Theorem 4.7 and Lemma 4.12]{CGHRangeI}.
The additional conditions for these to be regular-image factorizations
are given in \cite[Section 3.16 and Proposition 3.17]{CGHRangeII}.

The partial-map category $\Par(\mathcal D,\Mono)$ has morphisms given
by spans with monic left leg, composed by pullback
\cite[Section 3]{CockettLack}. Since $\mathcal D$ is regular, it is a
split restriction category with ranges
\cite[Theorem 4.5 and Section 4.11]{CGHRangeI}; its total subcategory
is canonically isomorphic to $\mathcal D$.
The sentence theory is represented in this category by
\[
 t_A=[S\xleftarrow{m_A}T_{A,S}\xrightarrow{m_A}S]:S\rightharpoonup S,
 \qquad m_A:T_{A,S}\hookrightarrow S.
\]
Every monomorphism in $\mathcal D$ has a partial inverse in this
partial-map category. Lemma~\ref{lem:range-regular-image} therefore
gives $t_A=\ran(c_{A,S})$.
If instead $\mathcal D=\Tot(\Split(\mathcal R))$ for a restriction
category $\mathcal R$ whose restriction splitting has ranges, require
$m_A$ to have a partial inverse
$r_A:S\rightharpoonup T_{A,S}$ in $\Split(\mathcal R)$.
The same lemma gives
\[
 \ran(c_{A,S})=m_A r_A=\bar r_A.
\]
Thus the range represents the constructed sentence theory.

A functor $G:\mathcal S\to\mathcal S'$ between restriction categories
with ranges is a \emph{restriction functor} if
$G(\bar v)=\overline{G(v)}$ for every morphism $v$.
If it also preserves ranges, then the splitting in
Lemma~\ref{lem:range-regular-image} satisfies
\[
 G(m)G(r)=\ran(G(f)),\qquad G(r)G(m)=1_{G(I)}.
\]
Hence $G(m)$ and $G(r)$ split the range of $G(f)$.

\subsection{Propositional theories}
\label{app:propositional-theories}

Let $\mathcal D$ be regular with finite coproducts, fix an atom
object $V\in\operatorname{Ob}(\mathcal D)$, and suppose
$G(Z)=V+Z+Z^2$ has an initial algebra
$S$ with constructors $\theta:V\to S$, $\nu:S\to S$, and
${\Rightarrow}:S^2\to S$. For an axiom subobject
$i_A:A\hookrightarrow S$, let $L=S^2+S^3+S^2$ and let
$\lambda:L\to S$ be given by
\eqref{eq:hilbert-one}--\eqref{eq:hilbert-three}, with $\nu$ for
negation. Use $R_0=L+A$, $q_0=[\lambda,i_A]$, $R_1=0$, and
\[
 R_2=S^2,\qquad
 p_2=\langle\pi_1,\Rightarrow\rangle:S^2\to S^2,
 \qquad q_2=\pi_2:S^2\to S.
\]
If the resulting inference endofunctor has an initial algebra,
Theorem~\ref{thm:least-theorem-subobject} constructs the least
subobject of $S$ containing the axioms and closed under these rules.
Its partial identity belongs to $\Par(\mathcal D,\Mono)$.
In particular, subobjects $i_+:A_+\hookrightarrow V$ and
$i_-:A_-\hookrightarrow V$ specify positive and negative atomic
axioms by the image of
$[\theta i_+,\nu\theta i_-]:A_++A_-\to S$.

\begin{thebibliography}{99}
\bibitem{Aczel}
P.~Aczel, \emph{An introduction to inductive definitions},
in J.~Barwise (ed.), Handbook of Mathematical Logic,
North-Holland, 1977, pp.~739--782.
\href{https://lawrencecpaulson.github.io/papers/Aczel-Inductive-Defs.pdf}{Text}.
\bibitem{CockettLack}
J.~R.~B. Cockett and S.~Lack, \emph{Restriction categories I:
Categories of partial maps}, Theoretical Computer Science
\textbf{270} (2002), pp.~223--259.
\href{https://doi.org/10.1016/S0304-3975(00)00382-0}{doi:10.1016/S0304-3975(00)00382-0}.
\bibitem{HP}
P.~H\'ajek and P.~Pudl\'ak, \emph{Metamathematics of First-Order
Arithmetic}, Perspectives in Logic, Cambridge University Press,
2016 reissue of the 1993 edition.
\href{https://doi.org/10.1017/9781316717271}{doi:10.1017/9781316717271}.
\bibitem{HH}
H.~P.~Hirst and J.~L.~Hirst, \emph{A Primer for Logic and Proof},
2015 edition.
\href{https://www.appstate.edu/~hirstjl/primer/hirst.pdf}{Authors' text}.
\bibitem{GTWW}
J.~A. Goguen, J.~W. Thatcher, E.~G. Wagner, and J.~B. Wright,
\emph{Initial algebra semantics and continuous algebras},
Journal of the ACM \textbf{24} (1977), no.~1, pp.~68--95.
\href{https://doi.org/10.1145/321992.321997}{doi:10.1145/321992.321997}.
\bibitem{GoguenBurstall}
J.~A. Goguen and R.~M. Burstall,
\emph{A study in the foundations of programming methodology:
Specifications, institutions, charters and parchments},
in D.~Pitt, S.~Abramsky, A.~Poign\'e, and D.~Rydeheard (eds.),
Category Theory and Computer Programming, Lecture Notes in Computer
Science \textbf{240}, Springer, 1986, pp.~313--333.
\href{https://doi.org/10.1007/3-540-17162-2_131}{doi:10.1007/3-540-17162-2\_131}.
\bibitem{CockettLists}
J.~R.~B. Cockett, \emph{List-arithmetic distributive categories: Locoi},
Journal of Pure and Applied Algebra \textbf{66} (1990), no.~1, pp.~1--29.
\href{https://doi.org/10.1016/0022-4049(90)90121-W}{doi:10.1016/0022-4049(90)90121-W}.
\bibitem{FPT}
M.~Fiore, G.~Plotkin, and D.~Turi,
\emph{Abstract syntax and variable binding},
Proceedings of the Fourteenth Annual IEEE Symposium on Logic in
Computer Science, IEEE Computer Society, 1999, pp.~193--202.
\href{https://doi.org/10.1109/LICS.1999.782615}{doi:10.1109/LICS.1999.782615}.
\bibitem{GabbayPitts}
M.~J. Gabbay and A.~M. Pitts,
\emph{A new approach to abstract syntax with variable binding},
Formal Aspects of Computing \textbf{13} (2002), pp.~341--363.
\href{https://doi.org/10.1007/s001650200016}{doi:10.1007/s001650200016}.
\bibitem{Katsumata}
S.-y. Katsumata,
\emph{A generalisation of pre-logical predicates to simply typed formal systems},
in J.~D\'iaz, J.~Karhum\"aki, A.~Lepist\"o, and D.~Sannella (eds.),
Automata, Languages and Programming, Lecture Notes in Computer Science
\textbf{3142}, Springer, 2004, pp.~831--845.
\href{https://doi.org/10.1007/978-3-540-27836-8_70}{doi:10.1007/978-3-540-27836-8\_70}.
\bibitem{MaiettiLists}
M.~E. Maietti, \emph{Joyal's arithmetic universe as list-arithmetic pretopos},
Theory and Applications of Categories \textbf{24} (2010),
no.~3, pp.~39--83.
\href{https://www.math.unipd.it/~maietti/papers/jartacrev.pdf}{Author's text}.
\bibitem{CGHRangeI}
J.~R.~B. Cockett, X.~Guo, and P.~Hofstra,
\emph{Range categories I: General theory},
Theory and Applications of Categories \textbf{26} (2012),
no.~17, pp.~412--452.
\href{https://www.tac.mta.ca/tac/volumes/26/17/26-17abs.html}{Journal page}.
\bibitem{CGHRangeII}
J.~R.~B. Cockett, X.~Guo, and P.~Hofstra,
\emph{Range categories II: Towards regularity},
Theory and Applications of Categories \textbf{26} (2012),
no.~18, pp.~453--500.
\href{https://www.tac.mta.ca/tac/volumes/26/18/26-18abs.html}{Journal page}.
\bibitem{Wissmann}
T.~Wi\ss{}mann, \emph{Minimality notions via factorization systems and examples},
Logical Methods in Computer Science \textbf{18} (2022),
no.~3, article~31, pp.~31:1--31:26.
\href{https://doi.org/10.46298/lmcs-18(3:31)2022}{doi:10.46298/lmcs-18(3:31)2022}.

\end{thebibliography}
\end{document}
